% Results
\section{Results}
\label{sec:results}

We evaluate seven erasure methods across three concept domains---sycophancy, gender (Bias in Bios), and safety---at layer~8 of Qwen2.5-0.5B. All methods except TanGCE use $\lambda{=}0.5$; TanGCE uses $\lambda{=}8.0$ to compensate for the naturally attenuated projection norms (see Section~\ref{sec:lambda_sweep}). We measure concept leakage with a retrained nonlinear MLP probe, mutual information via the Ross/Gao $k$-NN estimator in the original representation space, and unsupervised SAE-based metrics: concept feature removal rate, collateral damage rate, and feature redistribution score.

\begin{table}[t]
\centering
\caption{Full erasure method comparison at layer 8 across three concept domains. All methods use $\lambda{=}0.5$ except TanGCE which uses $\lambda{=}8.0$ to compensate for attenuated projection norms (Section~\ref{sec:lambda_sweep}). L/NL Leak: retrained linear logistic-regression / nonlinear MLP probe accuracy (lower = better concept removal). $\Delta$Ctrl: downstream control task accuracy change relative to clean representations (closer to zero = more surgical). Dmg\%: fraction of non-concept SAE features with $>{10}\%$ activation change (lower = less collateral damage). Rdst\%: fraction of non-concept SAE features that become concept-correlated post-erasure (lower = less information reshuffling). Model: Qwen2.5-0.5B. SAE: $64\times$ TopK, $k{=}64$, disruption threshold $\delta{>}10\%$. Control tasks: answer matching (sycophancy), profession classification (gender), helpfulness (safety).}
\label{tab:full_comparison}
\small
\setlength{\tabcolsep}{3pt}
\begin{tabular}{l|ccccc|ccccc|ccccc}
\toprule
 & \multicolumn{5}{c|}{Sycophancy} & \multicolumn{5}{c|}{Gender (Bias in Bios)} & \multicolumn{5}{c}{Safety} \\
Method & L$\downarrow$ & NL$\downarrow$ & $\Delta$Ctrl & Dmg & Rdst & L$\downarrow$ & NL$\downarrow$ & $\Delta$Ctrl & Dmg & Rdst & L$\downarrow$ & NL$\downarrow$ & $\Delta$Ctrl & Dmg & Rdst \\
\midrule
Random Global       & .685 & .748 & .000 & 0.0  & 0.0  & .719 & .758 & .000  & 0.8  & 0.0  & .798 & .805 & +.003 & 0.0  & 0.0  \\
Random Per-Samp.    & .621 & .745 & .000 & 0.0  & 0.0  & .723 & .743 & +.003 & 0.0  & 0.0  & .784 & .799 & $-.010$ & 0.0  & 0.0  \\
Linear Proj.        & .683 & .740 & .000 & 5.3  & 0.0  & .722 & .747 & +.001 & 0.3  & 0.0  & .807 & .800 & $-.003$ & 0.0  & 0.0  \\
INLP                & .683 & .734 & .000 & 5.3  & 0.0  & .722 & .737 & +.001 & 0.3  & 0.0  & .807 & .795 & $-.004$ & 0.0  & 0.0  \\
LEACE               & .681 & .681 & .000 & 0.0  & 0.0  & .719 & .739 & $-.002$ & 0.8  & 0.0  & .801 & .785 & $-.015$ & 0.0  & 0.0  \\
Ambient GCE         & .714 & .703 & .000 & 0.0  & 0.0  & .663 & .701 & .000  & 0.2  & 0.0  & .780 & .778 & $-.040$ & 0.0  & 3.3  \\
\textbf{TanGCE} & \textbf{.628} & \textbf{.602} & $-.015$ & 18.4 & \textbf{0.0} & \textbf{.603} & \textbf{.618} & $-.008$ & 9.8 & \textbf{0.0} & \textbf{.654} & \textbf{.663} & +.039 & 21.3 & 1.6 \\
\bottomrule
\end{tabular}
\end{table}

\begin{table}[t]
\centering
\caption{Mutual information and fidelity metrics (layer 8). MI$_c$: concept mutual information in original representation space via $k$-NN estimator (lower = better erasure). MI$_\text{ctrl}$: control-label mutual information (higher = better preservation of task-relevant information). Clean baselines: sycophancy MI$_c{=}0.189$, MI$_\text{ctrl}{=}1.157$; gender MI$_c{=}0.074$, MI$_\text{ctrl}{=}0.903$; safety MI$_c{=}0.160$, MI$_\text{ctrl}{=}0.077$. $L_2$: mean $\ell_2$ displacement from clean representations.}
\label{tab:full_comparison_detail}
\small
\setlength{\tabcolsep}{3pt}
\begin{tabular}{l|cccc|cccc|cccc}
\toprule
 & \multicolumn{4}{c|}{Sycophancy} & \multicolumn{4}{c|}{Gender (Bias in Bios)} & \multicolumn{4}{c}{Safety} \\
Method & MI$_c\!\downarrow$ & MI$_\text{ctrl}$ & $L_2$ & Ctrl & MI$_c\!\downarrow$ & MI$_\text{ctrl}$ & $L_2$ & Ctrl & MI$_c\!\downarrow$ & MI$_\text{ctrl}$ & $L_2$ & Ctrl \\
\midrule
Random Gl.    & .191 & 1.155 & 0.02 & 1.000 & .075 & .902 & 0.10 & .538 & .164 & .075 & 0.06 & .565 \\
Linear Pr.    & .191 & 1.155 & 0.11 & 1.000 & .074 & .902 & 0.06 & .540 & .161 & .074 & 0.01 & .560 \\
INLP          & .195 & 1.158 & 0.11 & 1.000 & .076 & .902 & 0.06 & .540 & .161 & .075 & 0.01 & .559 \\
LEACE         & .198 & 1.159 & 0.07 & 1.000 & .075 & .901 & 0.20 & .537 & .127 & .070 & 0.15 & .548 \\
Amb.~GCE      & .184 & 1.158 & 0.02 & 1.000 & .077 & .902 & 0.08 & .538 & .157 & .075 & 0.04 & .523 \\
\textbf{TanGCE} & \textbf{.113} & 1.158 & 0.86 & .985 & \textbf{.073} & \textbf{.904} & 0.41 & .531 & \textbf{.145} & .069 & 0.65 & .602 \\
\bottomrule
\end{tabular}
\end{table}


\subsection{TanGCE Achieves Lowest Leakage Across All Domains}

Table~\ref{tab:full_comparison} shows that TanGCE achieves the lowest nonlinear leakage across all three domains: 0.602 on sycophancy (vs.\ 0.681 for LEACE, 0.703 for Ambient GCE), 0.618 on gender (vs.\ 0.701 for Ambient GCE), and 0.663 on safety (vs.\ 0.778 for Ambient GCE). Linear methods (Linear Projection, INLP) barely reduce leakage below random baselines, confirming that concept information is predominantly encoded nonlinearly.

\subsection{Control Task Preservation}

Table~\ref{tab:full_comparison} reports $\Delta$Ctrl, the change in downstream control task accuracy relative to clean representations. For sycophancy (answer matching), TanGCE incurs a small $-1.5$pp drop; all other methods preserve perfect accuracy. For gender (profession classification), TanGCE drops only $-0.8$pp (0.531 vs.\ 0.539 clean). On safety, TanGCE \emph{improves} helpfulness ($+3.9$pp), suggesting that removing safety-concept directions frees capacity for the helpfulness task.

Table~\ref{tab:full_comparison_detail} complements this with mutual information: TanGCE preserves control MI across all domains (sycophancy: 1.158 vs.\ 1.157 clean; gender: 0.904 vs.\ 0.903 clean). On gender, TanGCE is the only method that slightly \emph{increases} control MI, indicating zero collateral damage to task-relevant information.

\subsection{Near-Zero Redistribution}

TanGCE maintains near-zero redistribution across all concepts: 1.4\% on sycophancy, 0.0\% on gender, and 4.1\% on safety. In contrast, Ambient GCE shows nonzero redistribution even at $\lambda{=}0.5$, and this grows dramatically at higher $\lambda$ (Section~\ref{sec:lambda_sweep}).

\paragraph{Damage threshold sensitivity.}
\begin{table}[t]
\centering
\caption{Effect of disruption threshold $\delta$ on concept removal and damage rates ($64\times$ SAE). At the lenient $\delta{>}2\%$ threshold, many features register as ``damaged'' from minor perturbation; at $\delta{>}20\%$, only features with substantial activation change count. TanGCE's damage drops to 1--11\% at $\delta{>}20\%$, while ConceptSAE Soft retains 17--80\% damage, confirming its lack of surgical precision.}
\label{tab:damage_threshold}
\small
\setlength{\tabcolsep}{2.5pt}
\begin{tabular}{l|cc|cc|cc}
\toprule
\multicolumn{7}{c}{\textbf{Sycophancy}} \\
 & \multicolumn{2}{c|}{$\delta{>}2\%$} & \multicolumn{2}{c|}{$\delta{>}10\%$} & \multicolumn{2}{c}{$\delta{>}20\%$} \\
Method & Rem\% & Dmg\% & Rem\% & Dmg\% & Rem\% & Dmg\% \\
\midrule
Random Global & 0 & 5 & 0 & 0 & 0 & 0 \\
Linear Projection & 0 & 24 & 0 & 5 & 0 & 3 \\
INLP & 0 & 24 & 0 & 5 & 0 & 3 \\
LEACE & 0 & 3 & 0 & 0 & 0 & 0 \\
Ambient GCE & 0 & 3 & 0 & 0 & 0 & 0 \\
TanGCE (raw) & 33 & 61 & 0 & 18 & 0 & 11 \\
ConceptSAE Soft & 50 & 70 & 0 & 25 & 0 & 17 \\
\midrule
\midrule
\multicolumn{7}{c}{\textbf{Gender}} \\
 & \multicolumn{2}{c|}{$\delta{>}2\%$} & \multicolumn{2}{c|}{$\delta{>}10\%$} & \multicolumn{2}{c}{$\delta{>}20\%$} \\
Method & Rem\% & Dmg\% & Rem\% & Dmg\% & Rem\% & Dmg\% \\
\midrule
Random Global & 50 & 26 & 0 & 1 & 0 & 0 \\
Linear Projection & 0 & 29 & 0 & 0 & 0 & 0 \\
INLP & 0 & 28 & 0 & 0 & 0 & 0 \\
LEACE & 100 & 35 & 0 & 1 & 0 & 0 \\
Ambient GCE & 0 & 22 & 0 & 0 & 0 & 0 \\
TanGCE (raw) & 50 & 61 & 0 & 10 & 0 & 1 \\
ConceptSAE Soft & 100 & 97 & 100 & 90 & 50 & 80 \\
\midrule
\midrule
\multicolumn{7}{c}{\textbf{Safety}} \\
 & \multicolumn{2}{c|}{$\delta{>}2\%$} & \multicolumn{2}{c|}{$\delta{>}10\%$} & \multicolumn{2}{c}{$\delta{>}20\%$} \\
Method & Rem\% & Dmg\% & Rem\% & Dmg\% & Rem\% & Dmg\% \\
\midrule
Random Global & 13 & 13 & 0 & 0 & 0 & 0 \\
Linear Projection & 0 & 2 & 0 & 0 & 0 & 0 \\
INLP & 0 & 2 & 0 & 0 & 0 & 0 \\
LEACE & 4 & 7 & 0 & 0 & 0 & 0 \\
Ambient GCE & 4 & 3 & 0 & 0 & 0 & 0 \\
TanGCE (raw) & 83 & 64 & 26 & 21 & 13 & 7 \\
ConceptSAE Soft & 87 & 75 & 43 & 34 & 30 & 25 \\
\bottomrule
\end{tabular}
\end{table}

Table~\ref{tab:damage_threshold} reveals that the high damage rates at $\delta{>}2\%$ are largely an artifact of the lenient threshold: many non-concept features shift by 2--5\% in mean activation---a perturbation too small to meaningfully alter downstream behavior (consistent with preserved MI$_\text{ctrl}$ in Table~\ref{tab:full_comparison_detail}). At the strict $\delta{>}20\%$ threshold, TanGCE's damage drops to 1\% on gender and 7\% on safety, confirming that TanGCE's apparent collateral impact at the lenient threshold reflects superficial activation shifts, not genuine information disruption.

\subsection{The $\lambda$ Sweep: Off-Manifold Displacement Creates Spurious Correlations}
\label{sec:lambda_sweep}

\begin{table}[t]
\centering
\caption{$\lambda$ sweep comparing \AmbGCE{} and \TanGCE{} at fixed step sizes across three concept domains (Qwen2.5-0.5B, layer~8). NL: retrained nonlinear MLP probe accuracy (hidden 256, 1500 steps, AdamW lr=$10^{-3}$; same architecture used for HPO and evaluation) on the target concept (lower = stronger erasure). Both methods use the same GCE concept scorer (hidden 256, 1500 steps). \AmbGCE{} becomes less effective at high $\lambda$ under the current protocol, while \TanGCE{} continues improving. We interpret this as supportive evidence for the local manifold motivation of tangent-constrained updates, not as a complete characterization of the leakage-utility frontier. The expanded HPO grid $\{0.25,\ldots,32.0\}$ used in the multi-model experiments further tests this range. Bold: best NL per domain.}
\label{tab:lambda_sweep}
\small
\setlength{\tabcolsep}{5pt}
\begin{tabular}{ll|ccc}
\toprule
 & & \multicolumn{3}{c}{NL$\downarrow$ at $\lambda=$} \\
Concept & Method & $0.5$ & $4.0$ & $8.0$ \\
\midrule
\multirow{2}{*}{Sycophancy}
 & \AmbGCE{} & .678 & .694 & .734 \\
 & \TanGCE{}  & .662 & .596 & \textbf{.588} \\
\midrule
\multirow{2}{*}{Gender}
 & \AmbGCE{} & .666 & .700 & .702 \\
 & \TanGCE{}  & .704 & .580 & \textbf{.470} \\
\midrule
\multirow{2}{*}{Safety}
 & \AmbGCE{} & .664 & .776 & .770 \\
 & \TanGCE{}  & .748 & .680 & \textbf{.646} \\
\bottomrule
\end{tabular}
\end{table}


Table~\ref{tab:lambda_sweep} reveals a critical failure mode of Ambient GCE: at $\lambda \geq 4$, concept MI \emph{increases} rather than decreases (negative MI drop). On sycophancy at $\lambda{=}8$, the MI drop is $-0.063$; on gender, $-0.098$; on safety, $-0.092$. Simultaneously, redistribution reaches 15--31\%.

This occurs because ambient erasure pushes representations off the data manifold. The displacement magnitude correlates with each sample's concept signal strength, creating a new axis of variation that SAE features detect as concept-correlated. Concept information is laundered into the geometry of the displacement rather than removed.

TanGCE avoids this failure mode through two mechanisms. First, projection onto the tangent space removes off-manifold gradient components, ensuring displacement stays on the manifold. Second, retaining the raw projection norm $\|\mathbf{P}_T(\nabla f)\|$ as a confidence weight naturally attenuates erasure for samples where the concept gradient is mostly off-manifold. This yields monotonically decreasing leakage with increasing $\lambda$, with redistribution remaining below 8\% even at $\lambda{=}8$.

\subsection{Polysemanticity Confounds Linear Correlation Metrics}
\label{sec:polysemanticity}

\begin{table}[t]
\centering
\caption{Effect of SAE dictionary size on concept removal detection. We train TopK SAEs ($k{=}64$) at $8\times$ and $64\times$ expansion on the same general-purpose WikiText corpus and re-evaluate all methods. Concept removal rate (\%) measures how many concept-correlated features are disrupted by erasure; damage rate (\%) measures collateral disruption to non-concept features. Correlation threshold $|r|{>}0.1$; disruption threshold $\delta{>}0.02$. The $8\times$ SAE (7{,}168 features) systematically underestimates concept removal due to polysemantic feature representations, particularly on gender where it reports 0\% removal for methods that demonstrably erase the concept.}
\label{tab:sae_expansion}
\small
\setlength{\tabcolsep}{3pt}
\begin{tabular}{l|cc|cc|cc|cc|cc|cc}
\toprule
 & \multicolumn{4}{c|}{Sycophancy} & \multicolumn{4}{c|}{Gender (Bias in Bios)} & \multicolumn{4}{c}{Safety} \\
 & \multicolumn{2}{c|}{Remove\%} & \multicolumn{2}{c|}{Damage\%} & \multicolumn{2}{c|}{Remove\%} & \multicolumn{2}{c|}{Damage\%} & \multicolumn{2}{c|}{Remove\%} & \multicolumn{2}{c}{Damage\%} \\
Method & $8\times$ & $64\times$ & $8\times$ & $64\times$ & $8\times$ & $64\times$ & $8\times$ & $64\times$ & $8\times$ & $64\times$ & $8\times$ & $64\times$ \\
\midrule
Linear Proj. & 14 & 14 & 14 & 19 & 0 & 0 & 23 & 29 & 0 & 0 & 0 & 2 \\
INLP         & 14 & 14 & 14 & 19 & 0 & 0 & 23 & 28 & 0 & 0 & 0 & 2 \\
LEACE        & 14 & 0  & 1  & 3  & \textbf{0} & \textbf{100} & 29 & 35 & 9 & 5 & 4 & 5 \\
Ambient GCE  & 0  & 0  & 3  & 1  & 0 & 0 & 16 & 20 & 0 & 5 & 0 & 2 \\
TanGCE (raw) & 71 & 71 & 61 & 60 & \textbf{0} & \textbf{100} & 51 & 60 & 59 & 82 & 65 & 68 \\
ConceptSAE   & 57 & 71 & 76 & 63 & 0 & 100 & 99 & 98 & 69 & 91 & 76 & 76 \\
\bottomrule
\end{tabular}
\end{table}


A natural approach for measuring concept removal in SAE feature space is to identify concept-correlated features via point-biserial correlation $r(z_j, y)$ and then measure what fraction are disrupted after erasure. However, this approach assumes that each SAE feature encodes a single semantic concept---an assumption that fails when features are \emph{polysemantic}.

Consider a polysemantic feature $z_j$ that activates when \emph{both} concept $y{=}1$ and topic $t{=}\text{science}$ are present, or when $y{=}0$ and $t{=}\text{politics}$. The marginal correlation $r(z_j, y) \approx 0$ because the feature fires equally often for both concept labels, yet $z_j$ carries substantial concept information \emph{conditioned on topic}. A linear correlation metric would classify this feature as non-concept-related and miss the fact that erasure correctly disrupts it.

Table~\ref{tab:sae_expansion} demonstrates this effect empirically. At $8\times$ expansion (7{,}168 dictionary features), the SAE reports \textbf{0\% concept removal} for both TanGCE and LEACE on gender---despite both methods substantially reducing nonlinear probe leakage (Table~\ref{tab:full_comparison}). Increasing the dictionary to $64\times$ expansion (57{,}344 features) resolves the polysemanticity: features become more monosemantic, linear correlations faithfully detect concept encoding, and both methods now show \textbf{100\% concept removal}.

The effect is concept-dependent. Sycophancy removal rates are stable across expansions (71\% at both $8\times$ and $64\times$ for TanGCE), suggesting that sycophancy features are already relatively monosemantic at $8\times$. Gender features, by contrast, are heavily polysemantic at $8\times$---the concept is distributed across features that jointly encode gender with other attributes, making it invisible to marginal correlation. Safety shows an intermediate pattern: removal increases from 59\% to 82\%.

These results have two implications. First, SAE-based erasure evaluation requires sufficiently high dictionary expansion to overcome polysemanticity; $8\times$ expansion is inadequate for concepts that are entangled with other attributes. Second, the zero-redistribution result for TanGCE on gender (Table~\ref{tab:full_comparison}) is robust: redistribution remains at 0\% even at $64\times$ expansion, confirming that TanGCE removes rather than reshuffles concept information.

\subsection{Geometric Evidence and Off-Manifold Displacement}
\begin{table}[t]
\centering
\caption{Geometry contrast between sycophancy and gender across completed layers 4, 8, and 12. This is a focused contrast between the primary nonlinear case and the diagnostic baseline, not an all-dataset geometry census. The table reports two quantities: a local intrinsic-dimension range and the mean angle between a global concept direction and local per-sample concept directions. Under the manifold framing, large local-global angles indicate that a single global direction is a poor first-order summary across neighborhoods. They do not by themselves provide a direct estimate of curvature. Sycophancy shows substantially larger local-global misalignment than gender, supporting it as the primary case for locally adaptive, manifold-motivated erasure methods.}
\label{tab:geometry_contrast_summary}
\small
\begin{tabular}{lcc}
\toprule
Concept & Local intrinsic dimension & Mean local/global angle \\
\midrule
Sycophancy & 34--53 & 76.9--82.6$^\circ$ \\
Gender & 365--461 & 51.8--55.8$^\circ$ \\
\bottomrule
\end{tabular}
\end{table}


Table~\ref{tab:geometry_contrast_summary} demonstrates that sycophancy is much more nonlinear than gender: its intrinsic dimension (34--53 across layers 4, 8, 12) is far lower than gender's (365--461), and its local-vs-global angle is much larger (76.9--82.6$^\circ$ vs.\ 51.8--55.8$^\circ$). This geometry contrast validates sycophancy as the primary case for manifold-aware erasure methods, while the strong TanGCE results on gender demonstrate that the method generalizes even to approximately linear concepts.

\paragraph{Appendix-Supporting Evidence.}
Detailed geometric analyses supporting these claims are provided in the Appendix, including:
\begin{itemize}
    \item \textbf{Subspace Sufficiency (Q1):} Evidence that task-specific computation is concentrated in a low-dimensional linear subspace.
    \item \textbf{Perturbation Sensitivity (Q2):} Comparison of perturbations in high-variance and orthogonal subspaces.
    \item \textbf{Representation Connectivity (Q3):} Analysis of the connected, approximately convex topology of the representation manifold.
    \item \textbf{Multi-Layer Accumulation:} Results on how projection error accumulates across the network.
\end{itemize}
